% Topic 2: Inverses and Determinants for 2x2 Matrices
% Part of the Matrix Fundamentals section

\subsection{Inverses and Determinants for $2 \times 2$ Matrices}

In standard algebra, the inverse of a number $x$ (where $x \neq 0$) is $\frac{1}{x}$, such that $x \times \frac{1}{x} = 1$. Similarly, in matrix algebra, the \textbf{inverse} of a square matrix $A$, denoted $A^{-1}$, is the matrix that multiplies with $A$ to yield the identity matrix $I$.

\[
A A^{-1} = A^{-1} A = I
\]

\subsubsection*{The Determinant}

The \textbf{determinant} of a matrix is a special scalar value that tells us whether the matrix has an inverse. For a $2 \times 2$ matrix:
\[
A = \begin{bmatrix} a & b \\ c & d \end{bmatrix}
\]
The determinant is denoted by $\det(A)$ or $|A|$ and is calculated as:
\[
\det(A) = ad - bc
\]

\subsubsection*{The Inverse Formula}

If $\det(A) \neq 0$, the inverse of the $2 \times 2$ matrix $A$ exists and is given by:
\[
A^{-1} = \frac{1}{\det(A)} \begin{bmatrix} d & -b \\ -c & a \end{bmatrix}
\]
Notice the pattern for finding the inverse of a $2 \times 2$ matrix:
\begin{itemize}[leftmargin=*]
  \item \textbf{Swap} the elements on the main diagonal ($a$ and $d$).
  \item \textbf{Negate} the elements on the other diagonal ($b$ and $c$).
  \item \textbf{Multiply} the resulting matrix by the scalar $\frac{1}{\det(A)}$.
\end{itemize}

\begin{remark}[Singular Matrices]
If $\det(A) = 0$, then $\frac{1}{\det(A)}$ is undefined. In this case, the matrix $A$ \textbf{does not have an inverse} and is called a \textbf{singular matrix}. If $\det(A) \neq 0$, the matrix is called \textbf{non-singular}.
\end{remark}

\bigskip

% ── Worked Example: Inverses and Determinants ──
\begin{problem}[Calculating Inverses and Determinants]
Consider the following two matrices:
\[
M = \begin{bmatrix} 4 & 1 \\ 2 & 3 \end{bmatrix}, \qquad N = \begin{bmatrix} 6 & -2 \\ -9 & 3 \end{bmatrix}
\]

\begin{enumerate}[label=\textbf{(\alph*)}]
  \item Calculate $\det(M)$. Hence, find $M^{-1}$.
  \item Verify that $M M^{-1} = I$.
  \item Determine whether matrix $N$ has an inverse.
\end{enumerate}
\end{problem}

\begin{solution}
\textbf{(a)} \\
First, find the determinant of $M$: \\
$\det(M) = (4)(3) - (1)(2) = 12 - 2 = 10$.

Since $\det(M) = 10 \neq 0$, the inverse exists. \\
Apply the $2 \times 2$ inverse formula (swap main diagonal, negate the other): \\
$M^{-1} = \frac{1}{10} \begin{bmatrix} 3 & -1 \\ -2 & 4 \end{bmatrix}$. \\
\emph{(Note: Leaving the scalar $1/10$ at the front is usually preferred unless decimals/fractions are specifically required for further calculation.)}

\medskip
\textbf{(b)} \\
To verify, calculate the product $M M^{-1}$: \\
$M M^{-1} = \begin{bmatrix} 4 & 1 \\ 2 & 3 \end{bmatrix} \left( \frac{1}{10} \begin{bmatrix} 3 & -1 \\ -2 & 4 \end{bmatrix} \right)$

Because scalar multiplication can be pulled to the front, we calculate the matrix product first: \\
$M M^{-1} = \frac{1}{10} \left( \begin{bmatrix} 4 & 1 \\ 2 & 3 \end{bmatrix} \begin{bmatrix} 3 & -1 \\ -2 & 4 \end{bmatrix} \right)$ \\[6pt]
$= \frac{1}{10} \begin{bmatrix} (4)(3)+(1)(-2) & (4)(-1)+(1)(4) \\ (2)(3)+(3)(-2) & (2)(-1)+(3)(4) \end{bmatrix}$ \\[6pt]
$= \frac{1}{10} \begin{bmatrix} 12-2 & -4+4 \\ 6-6 & -2+12 \end{bmatrix} = \frac{1}{10} \begin{bmatrix} 10 & 0 \\ 0 & 10 \end{bmatrix} = \begin{bmatrix} 1 & 0 \\ 0 & 1 \end{bmatrix} = I$. \checkmark

\medskip
\textbf{(c)} \\
First, find the determinant of $N$: \\
$\det(N) = (6)(3) - (-2)(-9) = 18 - 18 = 0$.

Because $\det(N) = 0$, matrix $N$ is a \textbf{singular matrix}. It does not have an inverse.
\end{solution}

\begin{takeaways}
\item \textbf{Determinant First:} Always calculate the determinant before attempting to find the inverse. If it is zero, you can state that the matrix is singular and stop immediately.
\item \textbf{Scalar Extraction:} When multiplying a matrix by its inverse to check your answer, it is much easier to leave the scalar fraction (like $\frac{1}{10}$) outside the matrix multiplication until the very end.
\item \textbf{Algebraic vs.\ Geometric:} A singular matrix ($\det = 0$) cannot be inverted algebraically (as this would require division by zero). Geometrically, this means the transformation squashes 2D space down into a 1D line or a point, losing information that cannot be uniquely reconstructed.
\end{takeaways}
